\chapter{Natural Disasters and Health Risks}

\medskip
\noindent\textit{\textbf{Under review.} This chapter is an AI-assisted educational draft; it is not a substitute for professional medical advice.}

\section*{Definition and Overview}
health effects of floods, fires, storms, heatwaves, earthquakes, and other natural hazards, including injury, infection, displacement, and psychological distress. Individual assessment is important because symptoms and treatment vary with the cause.

\section*{Clinical Features}
Injury, burns, smoke exposure, dehydration, heat illness, contaminated-water illness, worsening chronic disease, anxiety, and trauma symptoms.

\section*{When to Seek Medical Care}
Seek medical review for new, persistent, worsening, or function-limiting symptoms. Seek emergency help for severe breathing difficulty, collapse, sudden neurological change, severe chest or abdominal pain, or any rapidly worsening illness.

\section*{Diagnosis and Investigations}
Triage of immediate threats, assessment of injury and exposure, hydration and vital signs, infection assessment, and mental-health review when needed.

\section*{Management}
Follow emergency instructions, evacuate early when advised, obtain first aid, clean-water and medication supplies, and seek medical care for injury, breathing difficulty, heat illness, or severe distress.

\section*{Complications}
Possible complications depend on the cause and may include progression of the condition, effects on other organs, treatment adverse effects, and reduced quality of life. Early assessment can reduce preventable harm.

\section*{Prognosis and Outcomes}
Outcome depends on the underlying cause, severity, complications, and response to treatment. Discuss individual prognosis with the treating clinician.

\section*{Prevention and Self-Care}
Follow the treatment plan, keep relevant appointments, avoid known triggers, protect affected areas from injury, and seek help early if symptoms worsen. Do not start or stop prescription treatment without clinical advice.

\section*{Expanded clinical context}
Natural disasters is a prevention or health-behaviour topic. Interpretation should account for age, baseline health, medicines, comorbidities, goals, and access to care. This chapter supports discussion with a qualified clinician and does not establish a diagnosis.

\section*{Assessment and decision points}
Risk is shaped by environment, occupation, housing, access to care, habits, medicines, chronic disease, and social circumstances. Identify barriers and vulnerable people. Document the main concern, time course, functional impact, relevant risks, and red flags. Use targeted investigations when their result is likely to change management.

\section*{Management and follow-up}
Useful measures may include preparation, protective equipment, vaccination, safer routines, hydration, gradual activity, and support from health or community services. Explain expected improvement, when reassessment is needed, and how follow-up can be accessed. Avoid starting, stopping, or sharing prescription medicines without professional advice.

\section*{Safety-netting}
Seek urgent help for breathing difficulty, collapse, confusion, suspected poisoning, serious injury, or rapid deterioration after an exposure.

\section*{Practical prevention}
Review the plan regularly and use current local public-health and emergency guidance. Keep written instructions and seek review if the topic recurs, changes, or does not improve as expected.

\section*{Limitations}
This educational expansion should be checked against current local guidance and authoritative topic-specific sources.
\clearpage
